\SourceLeaf{283}%
\[
\begin{aligned}
l&=l^0+\lambda x+\lambda'xx+\text{etc.}\\
m&=m^0(1+\mu x+\mu'xx+\text{etc.})
\end{aligned}
\]
Then the calculation gives
\[
\begin{aligned}
\operatorname{tang}\psi={}&(1+\mu x+\mu'xx+\text{etc.})\operatorname{tang}\psi^0\\
&-l^0x-\tfrac{1}{2}(\lambda+l^0\mu)xx\\
&-(\tfrac{1}{3}\lambda'+\tfrac{1}{6}\lambda\mu-\tfrac{1}{6}l^0\mu\mu+\tfrac{2}{3}l^0\mu')x^3-\text{etc.}
\end{aligned}
\]
and from this, since \(u=\int\operatorname{tang}\psi\mathord{.}\mathrm{d}x\),
\[
\begin{aligned}
u={}&(x+\tfrac{1}{2}\mu xx+\tfrac{1}{3}\mu'x^3+\text{etc.})\operatorname{tang}\psi^0\\
&-\tfrac{1}{2}l^0xx-\tfrac{1}{6}(\lambda+l^0\mu)x^3\\
&-(\tfrac{1}{12}\lambda'+\tfrac{1}{24}\lambda\mu-\tfrac{1}{24}l^0\mu\mu+\tfrac{1}{6}l^0\mu')x^4-\text{etc.}
\end{aligned}
\]
where no constant need be added, because at \(x=0\) we also have \(u=0\). Since at \(x=h\) we likewise have \(u=0\), it follows from this equation that
\[
\operatorname{tang}\psi^0=\tfrac{1}{2}l^0h+(\tfrac{1}{6}\lambda-\tfrac{1}{12}l^0\mu)hh+(\tfrac{1}{12}\lambda'-\tfrac{1}{24}\lambda\mu)h^3+\text{etc.}
\]
Now take the equation for \(\psi\) and substitute for \(x\) the value \(h\), and for \(\operatorname{tang}\psi^0\) the value just found. This gives
\[
\begin{aligned}
\operatorname{tang}\psi'={}&-\tfrac{1}{2}l^0h-(\tfrac{1}{3}\lambda+\tfrac{1}{12}l^0\mu)hh\\
&-(\tfrac{1}{4}\lambda'+\tfrac{1}{24}\lambda\mu-\tfrac{1}{12}l^0\mu\mu+\tfrac{1}{6}l^0\mu')h^3\quad\text{etc.}
\end{aligned}
\]
\Needspace{19\baselineskip}
Since
\[
\begin{aligned}
l'&=l^0+\lambda h+\lambda'hh+\text{etc}\\
m'&=m^0(1+\mu h+\mu'hh+\text{etc.})
\end{aligned}
\]
we have
\[
\begin{aligned}
(\tfrac{1}{3}l^0+\tfrac{1}{6}l')h\BareRoot[6]{\frac{m^0}{m'}}={}&\tfrac{1}{2}l^0h+(\tfrac{1}{6}\lambda-\tfrac{1}{12}l^0\mu)hh\\
&+(\tfrac{1}{6}\lambda'-\tfrac{1}{36}\lambda\mu+\tfrac{7}{144}l^0\mu\mu-\tfrac{1}{12}l^0\mu')h^3\quad\text{etc.}
\end{aligned}
\]
\[
\begin{aligned}
-(\tfrac{1}{6}l^0+\tfrac{1}{3}l')h\BareRoot[6]{\frac{m'}{m^0}}={}&-\tfrac{1}{2}l^0h-(\tfrac{1}{3}\lambda+\tfrac{1}{12}l^0\mu)hh\\
&-(\tfrac{1}{3}\lambda'+\tfrac{1}{18}\lambda\mu-\tfrac{5}{144}l^0\mu\mu+\tfrac{1}{12}l^0\mu')h^3\quad\text{etc.}
\end{aligned}
\]
These thus agree in their first two terms, or through order \(hh\), with the above values of \(\operatorname{tang}\psi^0,\operatorname{tang}\psi'\). These convenient expressions can therefore be taken as sufficiently accurate values of the tangents, or, on including the factor \(206265''\), as the values of the angles \(\psi^0,\psi'\) themselves.

The length of the line \(L\) itself, between the points on the ellipsoid corresponding to the points \(F,G\) on the sphere, is the integral
